\documentclass[11.pt]{article}
\usepackage{amsmath, amssymb, amsfonts, amsthm}
\usepackage{pgfplots}
\usepackage{relsize}
\parindent 0px

\pgfplotsset{compat=1.18, width=12cm}

\begin{document}

\setcounter{section}{1}
\section{3D Perspective Projection}

Although we have a 2D coordinate representation of our window plane,\\
our points exist in 3 dimensions. Let's break things down and demonstrate\\
how to project a point in $\mathbb{R}^3$ onto a surface, converting it to
a point in $\mathbb{R}^2$.

Consider a rectangular pyramid oriented along the z-axis as shown in the\\
following plot.\\

\begin{tikzpicture}
\begin{axis}[
	title={Projection Pyramid},
	xlabel={$X$}, 
	ylabel={$Z$}, 
	zlabel={$Y$},
	xmin=-8,
	xmax=8,
	ymin=-1,
	ymax=5,
	zmin=-5,
	zmax=5,
	x axis line style={draw opacity=0},
	xtick=\empty,
	ytick={ 0,  1, 4},
	ztick={-4, -1, 0, 1, 4}
]
\draw[color=gray] (0,-1,0)--(0,5,0);
\draw[color=gray] (0,0,-5)--(0,0,5);

\draw[color=red] (0,0,0)--( 7.111, 4,  4);
\draw[color=red] (0,0,0)--(-7.111, 4,  4);
\draw[color=red] (0,0,0)--(-7.111, 4, -4);
\draw[color=red] (0,0,0)--( 7.111, 4, -4);

\draw[color=red] ( 7.111, 4,  4)--(-7.111, 4,  4);
\draw[color=red] (-7.111, 4,  4)--(-7.111, 4, -4);
\draw[color=red] (-7.111, 4, -4)--( 7.111, 4, -4);
\draw[color=red] ( 7.111, 4, -4)--( 7.111, 4,  4);

\draw[color=blue] ( 1.778, 1,  1)--(-1.778, 1,  1);
\draw[color=blue] (-1.778, 1,  1)--(-1.778, 1, -1);
\draw[color=blue] (-1.778, 1, -1)--( 1.778, 1, -1);
\draw[color=blue] ( 1.778, 1, -1)--( 1.778, 1,  1);

\end{axis}
\end{tikzpicture}

Our pyramid, shown in red, intercepts the z-axis at its apex and at the center\\
of its base. The rectangle in blue represents the screen of our program.\\

At the apex of the pyramid is a camera pointing at the application window.
The camera is within a pitch-black room, and it can only see the world 
through our program's 2D plane. The pyramid from beyond the window is the region\\
that the camera can see. However, the viewing region of the camera cuts off\\
abruptly at some distance from the window along the z-axis, forming the base\\
of the viewing pyramid.\\

Given this analogy, the apex of the figure represents the center of projection.\\
The blue rectangle is our view plane, where the 2D coordinates of our\\
application window lie. The base of the figure makes up the far plane. The\\
volume between the view plane and the far plane where we can see the outside\\
world is the viewing frustum.\\

The view plane and far plane can be treated as clipping planes, forming the\\
boundaries of our rendered world. However, as of the creation of this\\
document, I have not implemented a maximum z-coordinate value by which to\\
prevent further rendering. Currently, the graphics library only clips at the\\
viewing plane and renders infinitely beyond it.\\

\rule{12cm}{0.4pt}
\\

Now, we will look down the x-axis at the projection figure, such that we see\\
the yz-plane directly.\\

\begin{tikzpicture}
\begin{axis}[
	title={Pyramid Side-View},
	xlabel={$Z$},
	ylabel={$Y$},
	xmin=-1,
	xmax=5,
	ymin=-5,
	ymax=5,
	xtick={ 0, 1, 4},
	ytick={-4,-1, 0, 1, 4}
]
\draw [color=gray] (-1,0)--(5,0);
\draw [color=gray] (0,-5)--(0,5);

\draw[color=red] (0, 0)--(4, 4);
\draw[color=red] (0, 0)--(4,-4);
\draw[color=red] (4,-4)--(4, 4);

\draw[color=blue] (1,-1)--(1,1);

\addplot[only marks, color=gray] coordinates{(0,0)};
\addplot[only marks, color=gray] coordinates{(1,0)};
\addplot[only marks, color=gray] coordinates{(1,-1)};
\addplot[only marks, color=gray] coordinates{(1,1)};

%\draw[color=gray] (-0.5,-0.5) arc (-45:45:0.7071);
\draw[color=gray] (0.3536,-0.3536) arc (-45:45:0.5);

\node[below left]  at (axis cs: 0, 0) {$\vec{P_{proj}}$};
\node[below right] at (axis cs: 1, 0) {$\vec{P_{orig}}$};
\node[above left]  at (axis cs: 1, 1) {$\vec{P_{top}}$};
\node[below]  at (axis cs: 1,-1) {$\vec{P_{bottom}}$};

\end{axis}
\end{tikzpicture}

Using the plot above, we will define focal length $f$ and
field of view $FOV$.\\

Let $\vec{P_{proj}}=\begin{bmatrix} 0\\ 0\\ 0\end{bmatrix}$ be at the center
of projection.\\[5pt]

Let $\vec{P_{orig}}=\begin{bmatrix} 0\\ 0\\ f\end{bmatrix}$ be at the center of the
application window.\\[5pt]

Let $\vec{P_{bottom}}=\begin{bmatrix} 0\\ -1\\ f\end{bmatrix}$ be at the bottom of the
application window.\\[5pt]

Let $\vec{P_{top}}=\begin{bmatrix} 0\\ 1\\ f\end{bmatrix}$ be at the top of the
application window.\\[5pt]

The blue line segment $\overline{\vec{P_{bottom}}\vec{P_{top}}}$ represents the view
plane as seen directly from its side.\\

The focal length is the scalar distance between the center of projection and the
center of the application window, taken as the euclidean norm of the difference between 
both vectors.\\

$$ f = \left\lVert \vec{P_{orig}}-\vec{P_{proj}} \right\rVert$$\\

The field of view is the angle formed by points $\vec{P_{bottom}}$,
$\vec{P_{proj}}$, and $\vec{P_{top}}$.\\

$$ FOV = \mathlarger{\angle} \vec{P_{bottom}}\vec{P_{proj}}\vec{P_{top}}$$\\

Let us examine the relationship between $f$ and $FOV$.\\

$f$ is inversely proportional to $FOV$.\\

$$ f \propto \frac{1}{FOV}$$\\

The plot below demonstrates how the two properties interact.\\
The black dots represent 3 different centers of projection for\\
3 different projective regions. Each projective region has a\\ 
successively shorter focal length, and consequently, a successively\\ 
wider field of view. The size of the view plane for each region\\
remains constant.\\

\begin{tikzpicture}
\begin{axis}[
	title={$f$ and $FOV$ Relationship},
	xmin=0,
	xmax=4,
	ymin=-3,
	ymax=3,
	xtick=\empty,
	ytick=\empty
]
\draw (0,0)--(1,0);
\draw (1.5,0)--(2.5,0);
\draw (3,0)--(4,0);

\draw (0.75, -0.5)--(0.75, 0.5);
\draw (2.0, -0.5)--(2.0, 0.5);
\draw (3.25, -0.5)--(3.25, 0.5);

\addplot[only marks] coordinates{(0,0)};
\addplot[only marks] coordinates{(1.5,0)};
\addplot[only marks] coordinates{(3.0,0)};

\draw (0,0)--(1, 0.6667);
\draw (0,0)--(1, -0.6667);
\draw (1,-0.6667)--(1, 0.6667);

\draw (1.5, 0)--(2.5, 1);
\draw (1.5, 0)--(2.5, -1);
\draw (2.5, -1)--(2.5, 1);

\draw (3.0, 0)--(4.0, 2);
\draw (3.0, 0)--(4.0, -2);

\draw[dashed] (0, 0.5)--(4, 0.5);
\draw[dashed] (0, -0.5)--(4, -0.5);

\end{axis}
\end{tikzpicture}
\\

Let's return to our pyramid side view and calculate the focal length\\
from the field of view.\\

$$ \tan\left( \frac{1}{2}FOV \right)=
\frac{\left\lVert \vec{P_{top}}-\vec{P_{orig}}  \right\rVert}
     {\left\lVert \vec{P_{orig}}-\vec{P_{proj}} \right\rVert}$$\\

We can make 2 substitutions here. We already stated earlier that the\\
distance between the origin and the center of projection is the\\
focal length. In the previous section, we also stated that the distance\\
between the center and the top or bottom of our application window is 1.\\
Therefore,\\

$$ \left\lVert \vec{P_{top}}-\vec{P_{orig}}  \right\rVert = 1$$

Applying our substitutions,

$$ \tan\left( \frac{1}{2}FOV \right)=\frac{1}{f} $$

Solving for $f$,

$$ f = \frac{1}{\tan\left( \frac{1}{2}FOV \right)} $$\\

Finally, using the reciprocal trig identity for tangent,\\

$$ f = \cot\left( \frac{FOV}{2} \right) $$\\

Let's examine the boundaries for our FOV. We know that the FOV cannot be\\
less than 0 radians, as 0 would be the smallest possible FOV, which\\
is non-existent. We also know that FOV cannot be greater than $\tau$ radians,\\
because $\tau$ radians encircles the entire viewer. We can plot $f$ as a\\
function of $FOV$ from 0 to $\tau$ radians to investigate further.\\

\begin{tikzpicture}
\begin{axis}[
	title={$f(FOV)$ vs. $FOV$},
	xmin=0,
	xmax=2*pi,
	xtick={0, pi/2, pi, 3*pi/2, 2*pi},
	xticklabels={$0$ rad,$\frac{ \tau}{4}$ rad,
			     $\frac{ \tau}{2}$ rad,
			     $\frac{3\tau}{4}$ rad, $\tau$ rad},
	ymin=-4,
	ymax=4,
	ytick={-1, 0, 1}
]
\draw[gray] (0,-1)--(2*pi,-1);
\draw[gray] (0, 0)--(2*pi, 0);
\draw[gray] (0, 1)--(2*pi, 1);

\draw[gray] (pi/2  ,-4)--(pi/2  , 4);
\draw[gray] (pi    ,-4)--(pi    , 4);
\draw[gray] (3*pi/2,-4)--(3*pi/2, 4);

\addplot[domain=pi/32:2*pi-pi/32,
	 samples=300
] {cos(deg(x/2))/sin(deg(x/2))};

\end{axis}
\end{tikzpicture}

We need to consider the exclusivity of our boundaries.\\

$f$ has 2 vertical asymptotes at $FOV=0 \: rad$ and $FOV=\tau \: rad$. If we
express cotangent as the ratio cosine/sine, the focal-length can be written as\\

$$ f=\frac{ \cos\left( \frac{FOV}{2} \right) }
	  { \sin\left( \frac{FOV}{2} \right) } $$\\

At $FOV=0 \: rad$,

$$ \sin\left( \frac{FOV}{2} \right) = 
   \sin\left( \frac{ 0 }{2} \right) = 0, $$\\

which leaves $f$ undefined. At $FOV=\tau \: rad$,

$$ \sin\left( \frac{ FOV}{2} \right) = 
   \sin\left( \frac{\tau}{2} \right) = 0, $$\\

which also leaves $f$ undefined.\\

A naive analysis would end here, but we must evaluate where $f\leq0$. When\\
$f=0$, our projected coordinate values collapse to 0. When $f<0$, the signs\\
of our projected coordinates flip and the 2D image is mirrored both horizontally\\
and vertically. Using the cosine/sine expression listed above, we can evaluate\\
when the cosine function in the numerator falls below 0.\\

At $FOV=\frac{\tau}{2} \: rad$,

$$ \cos\left( \frac{ FOV}{2} \right) = 
   \cos\left( \frac{\tau}{4} \right) = 0, $$\\

which leaves $f$ at 0. For all values $\frac{\tau}{2}<FOV<\tau$,
$\cos\left( \frac{FOV}{2} \right) < 0$. Therefore,\\
these values can be excluded from our accepted FOV domain. The field of view\\
is limited to the following values:\\

$$\left\{ FOV\in\mathbb{R} \left| \: 0 \: rad <FOV<\frac{\tau}{2} \: rad \right. \right\}$$
\\

For testing purposes, it is convenient to use an FOV of $\frac{\tau}{4}$ radians. This
value gives us a focal length of 1, making our projective calculations fast and easy.\\

\rule{12cm}{0.4pt}
\\

Both X and Y coordinates are going to be projected from within the viewing\\ 
frustum onto the application screen. However, we can start things simple by\\
staying in our yz-plane and analyzing projections of the Y-coordinate alone.\\

\begin{tikzpicture}
\begin{axis}[
	title={Y-Coordinate Projections},
	xlabel={$Z$},
	ylabel={$Y$},
	xmin=-1,
	xmax=5,
	ymin=-5,
	ymax=5,
	xtick={ 0, 1, 1.5, 3.5, 4},
	ytick={-4,-1, 0, 1, 4}
]
\draw [color=gray] (-1,0)--(5,0);
\draw [color=gray] (0,-5)--(0,5);

\draw[color=red] (0, 0)--(4, 4);
\draw[color=red] (0, 0)--(4,-4);
\draw[color=red] (4,-4)--(4, 4);

\draw[color=blue] (1,-1)--(1,1);

\addplot[only marks, color=gray] coordinates{(0,0)};
\addplot[only marks, color=gray] coordinates{(1.5,1)};
\addplot[only marks, color=gray] coordinates{(3.5,1)};

\node[below left]  at (axis cs: 0, 0) {$\vec{P_{proj}}$};
\node[above right] at (axis cs: 1.5, 1) {$\vec{P_{z=1.5}}$};
\node[above] at (axis cs: 3.5, 1) {$\vec{P_{z=3.5}}$};

\draw[dashed, color=gray] (-1, 1)--(5, 1);
\draw[color=purple, opacity=0.5] (0,0)--(1.5,1);
\draw[color=purple, opacity=0.5] (0,0)--(3.5,1);

%\draw[color=gray] (-0.5,-0.5) arc (-45:45:0.7071);
\draw[color=gray] (0.3536,-0.3536) arc (-45:45:0.5);

\end{axis}
\end{tikzpicture}

Let $\vec{P_{z=1.5}}=\begin{bmatrix} 0\\ 1\\ 1.5 \end{bmatrix}$ be a point on the line
$\begin{bmatrix} X=0 & Y=1 \end{bmatrix}$

Let $\vec{P_{z=3.5}}=\begin{bmatrix} 0\\ 1\\ 3.5 \end{bmatrix}$ be a point on the same
line, further from the center of projection.\\

If we draw a straight line $\overline{\vec{P_{proj}}\vec{P_{z=1.5}}}$,
pictured in purple, then the projected coordinate values of $\vec{P_{z=1.5}}$ 
are at the intersection of this line and the view plane. The same goes for point
$\vec{P_{z=3.5}}$.\\

Note how projected Y-coordinate values get smaller as points get further from the
center of projection despite sharing the same Y-component values in 3-dimensional
space. If we continued extending points along $\begin{bmatrix} X=0 & Y=1 \end{bmatrix}$,
then the projected Y-component values would approach 0. The same can be said for the 
projected X-component values of points along $\begin{bmatrix} X=1 & Y=0 \end{bmatrix}$.
This means that we have a vanishing-point in $\mathbb{R}^2$ at 
$\vec{P_{2D}} = \begin{bmatrix} 0\\ 0\end{bmatrix}$ as Z-component values approach
infinity.

$$\lim_{z \to \infty} \vec{P_{2D}} = \begin{bmatrix} 0\\ 0\end{bmatrix}$$\\

The point $\vec{P_{orig}}$ on our view plane has coordinates 
$\begin{bmatrix} 0\\ 0\\ f\end{bmatrix}$ when evaluated in $\mathbb{R}^3$. However,
when evaluated in $\mathbb{R}^2$ we can ignore the Z-component. Using

$$\vec{P_{orig}} = \begin{bmatrix} 0\\ 0\end{bmatrix}$$\\

as our vanishing point, the limit can be written

$$\lim_{z \to \infty} \vec{P_{2D}} = \vec{P_{orig}}\text{.}$$\\

Now that we understand the mechanism behind 3D-coordinate projection, let's find
the value of the projected Y-Coordinate precisely.\\

\begin{tikzpicture}
\begin{axis}[
	title={View Plane Intercept},
	xlabel={$Z$},
	ylabel={$Y$},
	xmin=-1,
	xmax=5,
	ymin=-1,
	ymax=5,
	xtick={0, 1, 3.5},
	ytick={0, 0.42857, 1, 1.5},
	xticklabels={0, $f$, $Z_{3D}$},
	yticklabels={0, $Y_{2D}$, 1, $Y_{3D}$}
]
\draw [color=gray] (-1,0)--(5,0);
\draw [color=gray] (0,-1)--(0,5);

\draw[color=red] (0, 0)--(4, 4);
\draw[color=red] (4, 0)--(4, 4);
\draw[dashed, color=gray] (1, -1)--(1,0);
\draw[dashed, color=gray] (-1, 0.42857)--(1, 0.42857);
\draw[dashed, color=gray] (-1, 1.5)--(3.5,1.5);
\draw[dashed, color=gray] (3.5, -1)--(3.5,1.5);

\draw (0.80,0.00)--(0.80,0.20);
\draw (0.80,0.20)--(1.00,0.20);
\draw (3.30,0.00)--(3.30,0.20);
\draw (3.30,0.20)--(3.50,0.20);


\draw[color=blue] (1, 0)--(1,1);

\addplot[only marks, color=gray] coordinates{(0 , 0)};
\addplot[only marks, color=gray] coordinates{(1 , 0)};
\addplot[only marks, color=gray] coordinates{(3.5,1.5)};
\addplot[only marks, color=gray] coordinates{(1.0,0.42857)};
\addplot[only marks, color=gray] coordinates{(3.5,0.0)};
\node[below left]   at (axis cs: 0, 0) {$\vec{P_{proj}}$};
\node[below right]  at (axis cs: 0, 0) {$\alpha$};
\node[below right] at (axis cs: 1, 0) {$\vec{P_{orig}}$};
\node[right] at (axis cs: 3.5, 1.5) {$\vec{P_{3D}}$};
\node[below right] at (axis cs: 1.0, 0.42857) {$\vec{P_{2D}}$};
\node[below right] at (axis cs: 3.5, 0.0) {$\vec{P_{Z3D}}$};

\draw[color=purple, opacity=0.5] (0,0)--(3.5, 1.5);



%\draw[color=gray] (-0.5,-0.5) arc (-45:45:0.7071);
\draw[] (0.5, 0.0) arc (0:23.19859:0.5);

\end{axis}
\end{tikzpicture}

Let $\vec{P_{3D}} =  \begin{bmatrix}0\\Y_{3D}\\Z_{3D}\end{bmatrix}$ be a point within
the viewing frustum.\\[5pt]

Let $\vec{P_{2D}} =  \begin{bmatrix}0\\Y_{2D}\\f\end{bmatrix}$ be the 2-dimensional
projection of $\vec{P_{3D}}$.\\[5pt]

Let $\vec{P_{Z3D}} = \begin{bmatrix}0\\0\\Z_{3D}\end{bmatrix}$ be the point directly
below $\vec{P_{3D}}$ at the Z-axis.\\[5pt]

Let $\alpha = \mathlarger{\angle}\vec{P_{3D}}\vec{P_{proj}}\vec{P_{Z3D}}$ be the angle
formed at $\vec{P_{proj}}$ by the line $\overline{\vec{P_{proj}}\vec{P_{3D}}}$ and the
Z-axis.\\

Points $\vec{P_{3D}}$, $\vec{P_{proj}}$, and $\vec{P_{Z3D}}$ form a right triangle
that is parallel split by the view plane to produce another right triangle at points
$\vec{P_{2D}}$, $\vec{P_{proj}}$, and $\vec{P_{orig}}$.\\

$$ \triangle\vec{P_{3D}}\vec{P_{proj}}\vec{P_{Z3D}} \sim
   \triangle\vec{P_{2D}}\vec{P_{proj}}\vec{P_{orig}} $$\\

This pair of similar right triangles share the same angle $\alpha$ at the
center of projection. Therefore, we can derive the expression below:\\

$$ \tan(\alpha) = \frac{\left\lVert\vec{P_{3D}} - \vec{ P_{Z3D}}\right\rVert}
                       {\left\lVert\vec{P_{Z3D}}- \vec{P_{proj}}\right\rVert} = 
		  \frac{\left\lVert\vec{P_{2D}} - \vec{P_{orig}}\right\rVert}
		       {\left\lVert\vec{P_{orig}}-\vec{P_{proj}}\right\rVert}$$\\

Using the substitutions

$$\left\lVert\vec{P_{3D}} - \vec{ P_{Z3D}}\right\rVert = Y_{3D}$$
$$\left\lVert\vec{P_{Z3D}}- \vec{P_{proj}}\right\rVert = Z_{3D}$$
$$\left\lVert\vec{P_{2D}} - \vec{P_{orig}}\right\rVert = Y_{2D}$$
$$\left\lVert\vec{P_{orig}}-\vec{P_{proj}}\right\rVert = f$$\\

Our expression becomes\\

$$ \tan(\alpha) = \frac{Y_{3D}}{Z_{3D}} = \frac{Y_{2D}}{f}$$\\

We can disregard $\tan(\alpha)$, as it was merely a vehicle used to express a
relationship. Additionally, we can regard $Z_{3D} = Z$ as implicit. Using the simplest
form of our expression\\

$$ \frac{Y_{3D}}{Z} = \frac{Y_{2D}}{f} $$\\

we can solve for $Y_{2D}$\\

$$ Y_{2D} = \frac{Y_{3D}f}{Z}$$.\\

$X_{2D}$ can be calculated using a similar expression:

$$ X_{2D} = \frac{X_{3D}f}{Z}$$.\\

\rule{12cm}{0.4pt}
\\

There's one final element in our projective analysis that I haven't covered yet.\\

So far, we've treated FOV as a single angle formed at the center of projection 
by lines pointing to the edges of our application window. However, the application, 
or game window usually won't be a perfect square. How can we have 1 FOV for the X
and Y components of our view plane when the width and height of the screen are
different values? To put it simply, we can't. It is impossible to have a single angle
at the center of projection formed by lines that intercept the edges of the view plane
when the view plane is not square.\\

There are 3 separate ways that we can handle FOV. The following list outlines our 
approaches:

\begin{itemize}
\item Vertical FOV lock
\item Horizontal FOV lock
\item Anamorphic rendering
\end{itemize}

Each approach uses a different method to calculate the focal length. Consider 2
projective pyramids that share the same center of projection:

\begin{tikzpicture}
\begin{axis} [
	xmin = -1,
	xmax = 10,
	ymin = -1,
	ymax = 10,
	zmin = -1,
	zmax = 10,
	xtick = \empty,
	ytick = \empty,
	ztick = \empty
]

\end{axis}
\end{tikzpicture}

\end{document}
